% GENERATED FILE. Edit canonical lessons, then run npm run generate:books.
% canonical-source-hash: fnv1a64:19a631bcb35e8c76
% canonical-lessons: ES-C281-nuevo, ES-C281-viejo, ES-C281-limpio, ES-C281-sucio

\chapter{Four Ordinary Adjectives --- New, Old, Clean, Dirty}
\label{ch:cuatro-adjetivos}
\hlchaptermodality{281}

\begin{chapteropening}
\textbf{By the end of this chapter:} \emph{I can say new, old, clean and dirty in Spanish, and say what each of them meant in Latin before it meant what it means now.}

Complete ES-C281-sucio: say each adjective with its four endings, and give the Latin sense it started from.
\end{chapteropening}

\section[nuevo]{nuevo — the cousin test, one last time}

\label{lesson:ES-C281-nuevo}



\begin{quote}
Four adjectives to finish with, in two pairs. Each of them meant something slightly different in Latin from what it means to you now.
\end{quote}

\subsection*{What to know first}
\begin{itemize}
\raggedright
  \item grande — an adjective you already use.
  \item la puerta — where you first met stressed \emph{o} opening into \emph{ue}.
\end{itemize}

\begin{sounds}
\begin{itemize}
\raggedright
  \item \emph{NUE-vo}, stress on the first syllable, the \emph{ue} one glide. Four forms to keep straight: \emph{nuevo, nueva, nuevos, nuevas}. $\to$ reference
\end{itemize}
\end{sounds}

\begin{cousinweb}
\textbf{nuevo} — \textquotedblleft{}new.\textquotedblright{}

Latin \textbf{novus}, with the stressed \emph{o} opening into \emph{ue} exactly as it did in \emph{puerta}. Third time you have seen that change, and by now you should be able to run it backwards: give a Spanish \emph{ue}, guess a Latin \emph{o}.

English \textbf{new} is the cousin. It comes down the Germanic line from the same ancient root, never passing through Latin — which is why it is short, plain and worn, the way inherited words are.

The loans came later and stayed longer: a \textbf{novel} is a new thing, a \textbf{novice} is a new arrival, to \textbf{renovate} is to make new again, and a \textbf{nova} in the sky is a star that was not there last year. Every one of them is \emph{novus}, taken out of a book rather than out of a mouth.
\end{cousinweb}

\begin{grammarlens}[title={What you've built}]
You now have the cousin-and-loan test in a form you can apply without being told.

Short, irregular, everyday English words are usually the cousins. Long, tidy, bookish ones are usually the loans. \emph{New} against \emph{renovate}; \emph{wind} against \emph{ventilate}. When a Spanish word has two English relatives, the plain one tells you how old the family is and the fancy one tells you the Latin spelling.
\end{grammarlens}

\subsection*{Your turn}
Say these aloud:

\begin{itemize}
\raggedright
  \item \emph{nuevo, nueva, nuevos, nuevas}
  \item the Latin word, and the change that gave \emph{ue} (\emph{novus}; stressed \emph{o})
  \item \emph{una casa nueva}
\end{itemize}

\begin{tcolorbox}[breakable,colback=teal!4,colframe=teal!35!black,title={Before you move on}]
What Latin adjective is \emph{nuevo}? (\emph{Novus}.) Is English \emph{new} a loan from it? (No — a cousin.) Which English words are the loans? (\emph{Novel}, \emph{novice}, \emph{renovate}, \emph{nova}.)
\end{tcolorbox}
\section[viejo]{viejo — the word Spanish kept was the affectionate one}

\label{lesson:ES-C281-viejo}



\begin{quote}
The opposite of the word you just learned — and Spanish did not inherit the ordinary Latin adjective for it. It inherited a nickname.
\end{quote}

\subsection*{What to know first}
\begin{itemize}
\raggedright
  \item nuevo — its opposite, from \emph{novus}.
  \item trabajar — where you first met the throat sound written \emph{j}.
\end{itemize}

\begin{sounds}
\begin{itemize}
\raggedright
  \item \emph{VIE-jo}, stress on the first syllable. The \emph{j} is the same throat sound you have in \emph{trabajo}, not an English \emph{j}. $\to$ reference
\end{itemize}
\end{sounds}

\begin{cousinweb}
\textbf{viejo} — \textquotedblleft{}old.\textquotedblright{}

Latin's plain adjective was \textbf{vetus}. Spanish does not descend from it. It descends from \textbf{vetulus}, the diminutive — \emph{little old one}, the form you would use warmly about a grandparent or a well-worn coat. Spoken Latin shortened \emph{vetulus} to \emph{veclus}, and \emph{veclus} came out as \emph{viejo}, with the stressed \emph{e} opening into \emph{ie} along the way.

This happened more than once. Late spoken Latin liked diminutives the way a family likes nicknames, and several Spanish words are the affectionate form that outlived the neutral one. The plain \emph{vetus} did not vanish; it went into English instead, as \textbf{veteran} — someone grown old in a trade — and as \textbf{inveterate}, a habit that has aged into place.
\end{cousinweb}

\begin{grammarlens}[title={What you've built}]
Spanish inherited spoken Latin, not written Latin, and the two made different choices.

That is why the dictionary word so often ends up in English while the everyday word ends up in Spanish: the books preserved \emph{vetus}, the streets preserved \emph{vetulus}, and English later raided the books. When a Spanish word and its English relative feel different in tone, this is usually why.
\end{grammarlens}

\subsection*{Your turn}
Say these aloud:

\begin{itemize}
\raggedright
  \item \emph{viejo, vieja, viejos, viejas}
  \item \emph{nuevo} and \emph{viejo}, one after the other
  \item which Latin word English kept, and where (\emph{vetus}, in \emph{veteran})
\end{itemize}

\begin{tcolorbox}[breakable,colback=teal!4,colframe=teal!35!black,title={Before you move on}]
Which Latin form gave \emph{viejo} — the plain adjective or the diminutive? (The diminutive, \emph{vetulus}.) What did the diminutive add? (Warmth — \emph{little old one}.) Where did plain \emph{vetus} end up? (In English \emph{veteran}.)
\end{tcolorbox}
\section[limpio]{limpio — clear water, and what became of it}

\label{lesson:ES-C281-limpio}



\begin{quote}
The last pair. Both of these adjectives started out describing water, and only one of them stayed there.
\end{quote}

\subsection*{What to know first}
\begin{itemize}
\raggedright
  \item el agua — water, which this word first described.
  \item viejo — the adjective before this one.
\end{itemize}

\begin{sounds}
\begin{itemize}
\raggedright
  \item \emph{LIM-pio}, stress on the first syllable, and \emph{-pio} is one syllable, not two. $\to$ reference
\end{itemize}
\end{sounds}

\begin{cousinweb}
\textbf{limpio} — \textquotedblleft{}clean.\textquotedblright{}

Latin \textbf{limpidus}, which meant clear, transparent, see-through. It was said of water and of air, and it was about what you could see through rather than about what had been washed.

English borrowed it and left the meaning exactly where it found it: \textbf{limpid} is still a word for a clear pool or a clear prose style, and it has nothing to do with housework. Spanish, using the word every day rather than only in poetry, let it drift the short distance from \emph{clear} to \emph{clean}, and then handed the job of describing water to other words.

The verb followed the adjective: \emph{limpiar}, to clean, is built on \emph{limpio} and did not exist in Latin at all.
\end{cousinweb}

\begin{grammarlens}[title={What you've built}]
The same Latin word can sit still in one language and move in another, and the one that moved is usually the one that got used more.

There is a rough test in that. A word a language uses constantly drifts, narrows and widens; a word it reserves for books stays put. \emph{Limpid} is rare in English and has not changed in centuries. \emph{Limpio} is said a hundred times a day and has already travelled.
\end{grammarlens}

\subsection*{Your turn}
Say these aloud:

\begin{itemize}
\raggedright
  \item \emph{limpio, limpia, limpios, limpias}
  \item what \emph{limpidus} meant (clear, see-through)
  \item \emph{el agua limpia}
\end{itemize}

\begin{tcolorbox}[breakable,colback=teal!4,colframe=teal!35!black,title={Before you move on}]
What did \emph{limpidus} mean in Latin? (Clear, transparent.) Which English word still holds that sense? (\emph{Limpid}.) Which language moved the meaning, and why? (Spanish, because it used the word constantly.)
\end{tcolorbox}
\section[sucio]{sucio — wool straight off the sheep}

\label{lesson:ES-C281-sucio}



\begin{quote}
The opposite of the last word, and its history is far more specific than you would expect from something this ordinary.
\end{quote}

\subsection*{What to know first}
\begin{itemize}
\raggedright
  \item limpio — its opposite, from \emph{limpidus}.
  \item la mano — a body word, handy for practising this one.
\end{itemize}

\begin{sounds}
\begin{itemize}
\raggedright
  \item \emph{SU-cio}, stress on the first syllable. The \emph{c} before \emph{i} is a soft \emph{th} in much of Spain and an \emph{s} elsewhere, like the \emph{z} you met before. $\to$ reference
\end{itemize}
\end{sounds}

\begin{cousinweb}
\textbf{sucio} — \textquotedblleft{}dirty.\textquotedblright{}

Latin \textbf{sucidus}, from \textbf{sucus}, juice or sap. \emph{Sucidus} meant juicy, sappy, greasy — and its most ordinary use was for \textbf{wool that had not been washed}, still heavy with the natural grease it came off the animal with.

From there the road is short. Ungreased wool is unwashed wool; unwashed becomes unclean; unclean becomes dirty in general, including dirty in the moral sense that \emph{sucio} also carries. The word never described mud. It described a fleece.

English keeps \emph{sucus} in \textbf{succulent}, a plant full of juice, and the connection to \emph{suck} is real but ancient — both go back to a root about drawing liquid.
\end{cousinweb}

\begin{grammarlens}[title={What you've built}]
That is sixteen plain words, and every one of them turned out to have somewhere it came from.

Look back at what the pair you just finished has in common. \emph{Limpio} began at clear water and moved to clean; \emph{sucio} began at greasy wool and moved to dirty. Neither started life as an opposite of the other. They were pulled into a pair later, by speakers who needed two words and had these two lying about — which is how most of the neat oppositions in any language got built.
\end{grammarlens}

\subsection*{Your turn}
Say these aloud:

\begin{itemize}
\raggedright
  \item \emph{sucio, sucia, sucios, sucias}
  \item what \emph{sucidus} described (unwashed wool, still greasy)
  \item \emph{la mano sucia}, then \emph{la mano limpia}
\end{itemize}

\begin{tcolorbox}[breakable,colback=teal!4,colframe=teal!35!black,title={Before you move on}]
What did \emph{sucidus} mean? (Juicy or greasy — said of unwashed wool.) What is the noun behind it? (\emph{Sucus}, juice.) Were \emph{limpio} and \emph{sucio} ever a pair in Latin? (No — they were made into one later.)
\end{tcolorbox}
